\chapter{Álcoois: ativar o grupo OH}\label{ch:b1:alcohol-activation}

Os álcoois estão por toda parte: produzidos às toneladas a partir dos
alcenos e da fermentação, fornecidos por toda adição de Grignard do
capítulo anterior. No entanto, um álcool não sofre substituição como um
haloalcano: o grupo OH teria de sair como íon hidróxido, uma \omterm{def:b1:predominant-reaction:strong-acid}{base forte} e
um péssimo \omterm{def:b1:electronic-effects:nucleophile}{grupo de saída} (\cref{ch:b1:nucleophilic-substitution}).
Os químicos têm três respostas. Retirar o próton do OH, e o \omterm{def:b1:alcohol-activation:alkoxide}{alcóxido} se
torna um \omterm{def:b1:electronic-effects:nucleophile}{nucleófilo} poderoso. Adicionar um próton, e o \omterm{def:b1:electronic-effects:nucleophile}{grupo de saída}
passa a ser a água. Ou substituir o hidrogênio por um grupo que faz do
oxigênio parte de um excelente \omterm{def:b1:electronic-effects:nucleophile}{grupo de saída}, sem tocar no carbono. Este
capítulo segue os três caminhos, e as escolhas de mecanismo e de
estereoquímica que vêm com eles.

\begin{recall}
SN1, SN2 e sua estereoquímica (\cref{ch:b1:nucleophilic-substitution});
E1, E2 e a \omterm{prop:b1:elimination:zaitsev}{regra de Zaitsev} (\cref{ch:b1:elimination}); \omterm{def:b1:predominant-reaction:acidity-constant}{constantes de
acidez} e o método da \omterm{def:b1:predominant-reaction:predominant-reaction}{reação predominante} (\cref{ch:b1:predominant-reaction});
o efeito dos grupos alquila sobre a acidez dos álcoois
(\cref{exo:b1:electronic-effects:11}).
\end{recall}

\section{Os álcoois como ácidos: os alcóxidos}

\begin{definition}[Alcóxido]\label{def:b1:alcohol-activation:alkoxide}
Um \emph{alcóxido}\index{alcoxido@alcóxido} é a base conjugada \ce{RO-} de
um álcool \ce{ROH}: metóxido \ce{CH3O-}, etóxido \ce{CH3CH2O-},
terc-butóxido \ce{(CH3)3CO-}. Os alcóxidos são \omterm{def:b1:predominant-reaction:strong-acid}{bases fortes} e bons
\omterm{def:b1:electronic-effects:nucleophile}{nucleófilos}.
\end{definition}

\begin{proposition}[Acidez dos álcoois]\label{prop:b1:alcohol-activation:alcohol-acidity}
Os álcoois são ácidos mais fracos que a água, ou quase tão fracos quanto ela: $\mathrm{p}K_a$
15.5 (metanol), 15.9 (etanol), 17.1 (propan-2-ol), 19.2
(2-metilpropan-2-ol), contra 14.0 para o par \ce{H2O}/\ce{OH-}.
O hidróxido, portanto, só os desprotona em parte; uma conversão completa
em \omterm{def:b1:alcohol-activation:alkoxide}{alcóxido} exige uma base mais forte ou um metal alcalino.
\end{proposition}
% ledger: pka:methanol, pka:ethanol, pka:2-propanol, pka:tBuOH

\begin{proof}
A reação \ce{ROH + OH- <=> RO- + H2O} tem $K = 10^{14.0 -
\mathrm{p}K_a(\ce{ROH})}$, entre $10^{-1.5}$ e $10^{-5}$: desfavorável.
O sódio metálico reduz o próton, \ce{2ROH + 2Na -> 2RO- + 2Na+ + H2}, e
o íon hidreto do hidreto de sódio o captura de modo irreversível,
\ce{ROH + H- -> RO- + H2}: o hidrogênio escapa, e a reação é
completa qualquer que seja sua constante de equilíbrio.
\end{proof}

\begin{inthelab}[Hidreto de sódio]
O hidreto de sódio é vendido como um pó cinzento disperso em óleo mineral,
que o protege da umidade do ar. Ele reage violentamente com a água,
liberando hidrogênio inflamável; é pesado rapidamente, adicionado em porções
ao álcool seco sob gás inerte, e o borbulhamento do hidrogênio mostra o
avanço da desprotonação. O excesso de hidreto é destruído no fim pela
adição lenta de um álcool, nunca de água.
\end{inthelab}

\section{A síntese de Williamson de éteres}

\begin{definition}[Síntese de Williamson de éteres]\label{def:b1:alcohol-activation:williamson}
A \emph{síntese de Williamson de éteres}\index{sintese de Williamson de eteres@síntese de Williamson de éteres}
prepara um éter $\mathrm{R{-}O{-}R'}$ por uma reação SN2 de um \omterm{def:b1:alcohol-activation:alkoxide}{alcóxido} \ce{RO-}
com um haloalcano (ou um \omterm{def:b1:alcohol-activation:sulfonate}{éster sulfonato}) $\mathrm{R'{-}X}$.
\end{definition}

\begin{omfigure}
\setchemfig{atom sep=2.1em}
\schemestart
\chemfig{CH_3CH_2-@{o}O^{-}}\,\chemfig{Na^{+}}
\+
\chemfig{@{c}CH_3-[@{ci}]@{i}I}
\arrow{->}
\chemfig{CH_3CH_2-O-CH_3}
\+
\chemfig{Na^{+}}\,\chemfig{I^{-}}
\schemestop

\medskip
\chemmove{\draw[->, thick, shorten <=2pt, shorten >=2pt](o) .. controls +(80:8mm) and +(100:8mm) .. (c);
          \draw[->, thick, shorten <=2pt, shorten >=2pt](ci) .. controls +(90:5mm) and +(90:5mm) .. (i);}

{\small Síntese de Williamson do metoxietano: o etóxido ataca o carbono
do iodometano pelo lado oposto ao iodo (SN2).}
\end{omfigure}

\begin{method}[Escolher os dois parceiros]\label{met:b1:alcohol-activation:williamson-choice}
Um éter $\mathrm{R{-}O{-}R'}$ pode ser cortado de um lado ou do outro do oxigênio.
\begin{enumerate}
  \item Escreva os dois pares possíveis: \ce{RO-} com $\mathrm{R'X}$, e
        $\mathrm{R'O^-}$ com \ce{RX}.
  \item Fique com o par em que o haleto é metílico ou primário: a SN2
        precisa de um carbono pouco impedido.
  \item Rejeite um par com haleto terciário: o \omterm{def:b1:alcohol-activation:alkoxide}{alcóxido}, uma \omterm{def:b1:predominant-reaction:strong-acid}{base forte},
        provocaria eliminação (E2).
  \item Com um haleto secundário, espere uma mistura de éter e alceno.
\end{enumerate}
\end{method}

\begin{example}[Um éter assimétrico]\label{ex:b1:alcohol-activation:williamson}
2-Metoxi-2-metilpropano (éter metil-terc-butílico): o terc-butóxido de
sódio com o iodometano o forma de modo limpo; o metóxido de sódio com o
brometo terciário, o 2-bromo-2-metil\-propano, dá 2-metilpropeno por E2.
Volumoso do lado do \omterm{def:b1:alcohol-activation:alkoxide}{alcóxido}, pequeno do lado do haleto.
\end{example}

\section{Ativar o grupo OH}

\begin{definition}[Ésteres sulfonatos]\label{def:b1:alcohol-activation:sulfonate}
Um \emph{éster sulfonato}\index{ester sulfonato@éster sulfonato} $\mathrm{R{-}O{-}SO_2{-}R'}$ forma-se
a partir de um álcool e de um cloreto de sulfonila $\mathrm{R'SO_2Cl}$ na presença de
uma base (piridina). O \emph{tosilato}\index{tosilato}, \ce{R-OTs}, vem
do cloreto de 4-metilbenzenossulfonila (cloreto de tosila); o
\emph{mesilato}\index{mesilato}, \ce{R-OMs}, do cloreto de metanossulfonila.
O ânion sulfonato $\mathrm{R'SO_3^-}$, base conjugada de um \omterm{def:b1:predominant-reaction:strong-acid}{ácido forte} e
estabilizado por ressonância sobre três oxigênios, é um excelente \omterm{def:b1:electronic-effects:nucleophile}{grupo de saída}.
\end{definition}

\begin{omfigure}
\setchemfig{atom sep=2em}
\schemestart
\chemfig{R-O-H}
\+
\chemfig{Cl-S(=[2]O)(=[6]O)-[:0]*6(-=-(-CH_3)=-=)}
\arrow{->[piridina]}
\chemfig{R-O-S(=[2]O)(=[6]O)-[:0]*6(-=-(-CH_3)=-=)}
\schemestop
\par\vspace{6pt}

{\small Tosilação de um álcool. A ligação O--H é substituída por O--S; a
ligação C--O do álcool não é tocada. A piridina captura o HCl formado.}
\end{omfigure}

\begin{proposition}[A configuração na tosilação e na substituição]\label{prop:b1:alcohol-activation:tosylation-retention}
A tosilação de um álcool em um carbono estereogênico conserva sua
\omterm{def:b1:stereochemistry-in-depth:isomers}{configuração} (retenção); uma substituição SN2 posterior do \omterm{def:b1:alcohol-activation:sulfonate}{tosilato} a
inverte. As duas etapas juntas substituem OH por um \omterm{def:b1:electronic-effects:nucleophile}{nucleófilo} com inversão.
\end{proposition}

\begin{proof}
Na tosilação, as ligações que mudam são O--H e S--Cl: o oxigênio do
álcool ataca o enxofre e seu hidrogênio vai para a base. Nenhuma ligação
do carbono estereogênico é rompida, de modo que o arranjo em torno dele
permanece o mesmo. Na etapa SN2, a ligação C--O se rompe enquanto o
\omterm{def:b1:electronic-effects:nucleophile}{nucleófilo} se liga pelo lado oposto (\cref{prop:b1:nucleophilic-substitution:sn2-inversion}).
\end{proof}

\begin{omfigure}
\begin{tikzpicture}[font=\small, >={Stealth}]
  \node (a) at (0,0) {\chemfig{C(-[:150]H_3C)(<[:210]H)(<:[:250]C_2H_5)-OH}};
  \node at (0,-1.4) {\cip{S}-butan-2-ol};
  \draw[->, thick] (1.3,0) -- (2.5,0) node[midway, above] {TsCl} node[midway, below, font=\scriptsize] {retenção};
  \node (b) at (3.9,0) {\chemfig{C(-[:150]H_3C)(<[:210]H)(<:[:250]C_2H_5)-OTs}};
  \node at (3.9,-1.4) {tosilato \cip{S}};
  \draw[->, thick] (5.3,0) -- (6.6,0) node[midway, above] {\ce{CN-}} node[midway, below, font=\scriptsize] {SN2, inversão};
  \node (c) at (8.3,0) {\chemfig{[:0]NC-C(-[:30]CH_3)(<[:-30]H)(<:[:-70]C_2H_5)}};
  \node at (8.3,-1.4) {nitrila \cip{R}};
\end{tikzpicture}

{\small Do \cip{S}-butan-2-ol à 2-metilbutanonitrila: a tosilação
conserva a \omterm{def:b1:stereochemistry-in-depth:isomers}{configuração}, a etapa SN2 com o cianeto a inverte. No total:
OH substituído por CN com inversão.}
\end{omfigure}

O mesmo objetivo pode ser atingido em meio ácido: o grupo OH, protonado,
sai como água, \ce{R-OH2+ -> R+ + H2O} (SN1, E1) ou sob ataque (SN2). O
ácido, porém, limita o \omterm{def:b1:electronic-effects:nucleophile}{nucleófilo} aos que sobrevivem a ele --- íons
haleto, água, álcoois --- e abre a porta aos \omterm{def:b1:electronic-effects:intermediates}{carbocátions} e às suas
reações secundárias.

\section{Conversão em haloalcanos}

\begin{proposition}[Álcoois e haletos de hidrogênio]\label{prop:b1:alcohol-activation:hx-by-class}
Os álcoois terciários dão o haloalcano com ácido clorídrico concentrado à
temperatura ambiente (SN1); os álcoois primários precisam de ácido
bromídrico ou iodídrico e de aquecimento (SN2 sobre o álcool protonado);
os álcoois secundários reagem com velocidades intermediárias, muitas
vezes pelos dois mecanismos.
\end{proposition}

\begin{proof}
A protonação transforma OH em \ce{-OH2+}. Para um álcool terciário, a
água sai facilmente, dando um \omterm{def:b1:electronic-effects:intermediates}{carbocátion} estável capturado pelo haleto
(SN1): até o modesto \omterm{def:b1:electronic-effects:nucleophile}{nucleófilo} \ce{Cl-} basta. Um \omterm{def:b1:electronic-effects:intermediates}{carbocátion} primário é
instável demais; o haleto precisa empurrar a água para fora por trás
(SN2), o que exige um bom \omterm{def:b1:electronic-effects:nucleophile}{nucleófilo} (\ce{Br-}, \ce{I-}, melhores que
\ce{Cl-} em meio prótico,
\cref{prop:b1:electronic-effects:nucleophilicity-basicity}) e calor.
\end{proof}

\begin{example}[Um teste de classificação]\label{ex:b1:alcohol-activation:lucas}
Uma mistura de ácido clorídrico concentrado e cloreto de zinco (um \omterm{def:b1:lewis-resonance-vsepr:lewis-acid}{ácido
de Lewis} que ajuda o OH a sair) turva um álcool terciário de imediato ---
o cloreto insolúvel se separa ---, um álcool secundário em alguns minutos,
e um álcool primário quase nada à temperatura ambiente: a ordem de
estabilidade dos \omterm{def:b1:electronic-effects:intermediates}{carbocátions} tornada visível em um tubo de ensaio.
\end{example}

\section{Desidratação e formação de éteres em meio ácido}

\begin{definition}[Desidratação de um álcool]\label{def:b1:alcohol-activation:dehydration}
A \emph{desidratação}\index{desidratacao@desidratação} de um álcool é a
eliminação de água, catalisada por um \omterm{def:b1:predominant-reaction:strong-acid}{ácido forte} (ácido sulfúrico ou
fosfórico) e pelo calor, que dá um alceno: $\mathrm{R_2CH{-}CR'_2OH} \to \mathrm{R_2C{=}CR'_2} + \ce{H2O}$.
\end{definition}

\begin{omfigure}
\setchemfig{atom sep=2.1em}
\schemestart
\chemfig{C(-[2]CH_3)(-[:210]H_3C)(-[:-30]CH_2CH_3)-OH}
\arrow{->[\ce{H+}]}
\chemfig{C(-[2]CH_3)(-[:210]H_3C)(-[:-30]CH_2CH_3)-OH_2^{+}}
\arrow{->[\ce{-H2O}]}
\chemfig{C^{+}(-[2]CH_3)(-[:210]H_3C)(-[:-30]CH_2CH_3)}
\schemestop

\medskip
\setchemfig{arrow coeff=1.8}
\schemestart
\chemfig{C^{+}(-[2]CH_3)(-[:210]H_3C)(-[:-30]CH_2CH_3)}
\arrow{->[\ce{-H+}][(majoritário)]}
\chemfig{C(-[2]CH_3)(-[:210]H_3C)=[:-30]CHCH_3}
\schemestop

\medskip
{\small \omterm{def:b1:alcohol-activation:dehydration}{Desidratação} do 2-metilbutan-2-ol catalisada por ácido (E1):
protonação, perda de água até o \omterm{def:b1:electronic-effects:intermediates}{carbocátion} terciário, perda de um próton.
Retirar o próton do \ce{CH2} dá o 2-metilbut-2-eno trissubstituído
(Zaitsev, majoritário); retirá-lo de um grupo metila, o 2-metilbut-1-eno
(minoritário).}
\end{omfigure}

\begin{proposition}[Éter ou alceno]\label{prop:b1:alcohol-activation:ether-vs-alkene}
Um álcool primário aquecido com ácido sulfúrico dá principalmente um éter
simétrico a temperatura moderada (uma molécula substituindo outra por
SN2) e principalmente um alceno a temperatura mais alta; os álcoois
secundários e terciários dão alcenos.
\end{proposition}

\begin{proof}
Para um álcool primário, o álcool protonado pode ser atacado no carbono
por uma segunda molécula de álcool (SN2, dando \ce{R-O-R} após a perda de
um próton) ou em um hidrogênio $\beta$ (E2). A eliminação produz mais
moléculas do que consome e é favorecida à medida que a temperatura sobe
(\cref{prop:b1:elimination:sn-vs-e}); a substituição domina quando ela é
mais baixa. Os álcoois secundários e terciários se ionizam em
\omterm{def:b1:electronic-effects:intermediates}{carbocátions}, que perdem um próton mais facilmente do que são capturados
por um \omterm{def:b1:electronic-effects:nucleophile}{nucleófilo} fraco, tanto mais quanto mais se aquece; a água formada
também é removida pela destilação do alceno volátil.
\end{proof}

\section{Exercícios}

\begin{exercise}[$\star$]\label{exo:b1:alcohol-activation:1}
Escreva as reações do etanol com o sódio, com o hidreto de sódio e com o
hidróxido de sódio; calcule a constante da última ($\mathrm{p}K_a$ do
etanol 15.9).
\end{exercise}

\begin{exercise}[$\star$]\label{exo:b1:alcohol-activation:2}
Dê os produtos de: metóxido de sódio com 1-bromopropano; etóxido de
sódio com iodometano; fenóxido de sódio com bromoetano.
\end{exercise}

\begin{exercise}[$\star$]\label{exo:b1:alcohol-activation:3}
Escreva a tosilação do propan-1-ol e a reação do \omterm{def:b1:alcohol-activation:sulfonate}{tosilato} com o iodeto
de sódio. Qual é a vantagem sobre a reação direta do álcool com o iodeto
de sódio?
\end{exercise}

\begin{exercise}[$\star$]\label{exo:b1:alcohol-activation:4}
Dê os alcenos formados pela \omterm{def:b1:alcohol-activation:dehydration}{desidratação} ácida do butan-2-ol e do
2-metilpentan-2-ol, e o majoritário.
\end{exercise}

\begin{exercise}[$\star\star$]\label{exo:b1:alcohol-activation:5}
Proponha uma síntese de Williamson do 1-etoxi-2-metilpropano e explique
sua escolha de parceiros. Por que a outra escolha daria um alceno?
\end{exercise}

\begin{exercise}[$\star\star$]\label{exo:b1:alcohol-activation:6}
O \cip{R}-octan-2-ol é convertido em seu \omterm{def:b1:alcohol-activation:sulfonate}{tosilato} e depois tratado com
etóxido de sódio em etanol. Dê a \omterm{def:b1:stereochemistry-in-depth:isomers}{configuração} do éter. O que se
obteria aquecendo o mesmo álcool com ácido sulfúrico?
\end{exercise}

\begin{exercise}[$\star\star$]\label{exo:b1:alcohol-activation:7}
Escreva o mecanismo da reação do 2-metilpropan-2-ol com ácido clorídrico
concentrado e explique por que a reação é rápida à temperatura
ambiente.
\end{exercise}

\begin{exercise}[$\star\star$]\label{exo:b1:alcohol-activation:8}
Escreva o mecanismo da formação do etoxietano a partir do etanol em
ácido sulfúrico, e o da eliminação concorrente.
\end{exercise}

\begin{exercise}[$\star\star$]\label{exo:b1:alcohol-activation:9}
No teste de classificação deste capítulo, qual álcool, entre o
butan-1-ol, o butan-2-ol e o 2-metilpropan-2-ol, turva primeiro? Escreva
o produto.
\end{exercise}

\begin{exercise}[$\star\star\star$]\label{exo:b1:alcohol-activation:10}
O 3,3-dimetilbutan-2-ol, aquecido em meio ácido, dá principalmente o
2,3-dimetilbut-2-eno, cujo esqueleto carbônico difere do do álcool.
Proponha um mecanismo com um deslocamento 1,2 de um grupo metila no
\omterm{def:b1:electronic-effects:intermediates}{carbocátion}, e justifique esse deslocamento.
\end{exercise}

\begin{exercise}[$\star\star\star$]\label{exo:b1:alcohol-activation:11}
Planeje uma síntese do \cip{S}-2-metoxibutano a partir do
\cip{R}-butan-2-ol, e outra a partir do \cip{S}-butan-2-ol, mantendo em
cada uma a \omterm{def:b1:stereochemistry-in-depth:isomers}{configuração} sob controle.
\end{exercise}

\begin{exercise}[$\star\star\star$]\label{exo:b1:alcohol-activation:12}
Mostre que o equilíbrio \ce{2CH3CH2OH <=> (CH3CH2)2O + H2O} pode ser
deslocado para o éter removendo a água continuamente, usando o
\omterm{def:b1:extent-q-and-k:quotient}{quociente de reação} (\cref{ch:b1:extent-q-and-k}).
\end{exercise}

\section{Problema: dois caminhos para um éter}

\begin{problem}[{Problema de fim de semana --- as duas rotas de Williamson
para o éter metil-terc-butílico, um éter de estereoquímica controlada a
partir de um álcool opticamente ativo, a desidratação de um álcool
terciário e a massa de éter obtida}]
\label{pb:b1:alcohol-activation:1}
O 2-metoxi-2-metilpropano (éter metil-terc-butílico, MTBE) foi produzido
em grande escala como aditivo de combustível. Massas molares (\unit{g/mol}): \ce{C4H10O}
74.0, \ce{C5H12O} 88.0, \ce{CH3I} 141.9.

\medskip
\noindent\textbf{Parte I --- O MTBE pela síntese de Williamson.}
\begin{enumerate}
  \item Escreva os dois pares de parceiros (\omterm{def:b1:alcohol-activation:alkoxide}{alcóxido} e haleto) que
        poderiam dar o MTBE.
  \item Qual par fracassa? Escreva a reação que ocorre em seu lugar.
  \item Escreva a reação bem-sucedida e seu mecanismo.
  \item Como se prepara o terc-butóxido de sódio a partir do
        2-metilpropan-2-ol? Por que não com hidróxido de sódio?
  \item Por que a reação é feita no álcool seco ou em um \omterm{def:b1:intermolecular-forces-solvents:solvent-classes}{solvente aprótico}?
  \item O produto é \omterm{def:b1:stereochemistry-in-depth:stereocentre}{quiral}?
\end{enumerate}

\medskip
\noindent\textbf{Parte II --- Um éter opticamente ativo.}
\begin{enumerate}[resume]
  \item O \cip{R}-butan-2-ol é convertido em seu \omterm{def:b1:alcohol-activation:sulfonate}{tosilato}. Escreva a
        reação e dê a \omterm{def:b1:stereochemistry-in-depth:isomers}{configuração} do \omterm{def:b1:alcohol-activation:sulfonate}{tosilato}.
  \item O \omterm{def:b1:alcohol-activation:sulfonate}{tosilato} reage com o metóxido de sódio. Escreva o produto e
        sua \omterm{def:b1:stereochemistry-in-depth:isomers}{configuração}.
  \item Qual \omterm{def:b1:stereochemistry-in-depth:isomers}{configuração} do álcool dá o \cip{R}-2-metoxibutano pelo
        mesmo caminho?
  \item Outro caminho: o \omterm{def:b1:alcohol-activation:alkoxide}{alcóxido} do \cip{R}-butan-2-ol com o
        iodometano. \omterm{def:b1:stereochemistry-in-depth:isomers}{Configuração} do produto? Por quê?
  \item Que reação secundária concorre na questão 8, e por que ela é
        limitada aqui?
  \item Por que aquecer o \cip{R}-butan-2-ol com metanol em ácido
        sulfúrico daria um éter quase racêmico, se é que daria algum?
\end{enumerate}

\medskip
\noindent\textbf{Parte III --- \omterm{def:b1:alcohol-activation:dehydration}{Desidratação} de um álcool terciário.}
\begin{enumerate}[resume]
  \item Escreva o mecanismo da \omterm{def:b1:alcohol-activation:dehydration}{desidratação} ácida do
        2-metilbutan-2-ol.
  \item Dê os dois alcenos e preveja o majoritário.
  \item Por que nenhum éter se forma a partir deste álcool?
  \item Por que o 2-metilbut-2-eno é destilado à medida que se forma?
  \item Desenhe o \omterm{def:b1:elementary-steps:energy-profile}{perfil de energia} da etapa E1 e indique a \omterm{def:b1:elementary-steps:rds}{etapa
        determinante da velocidade}.
  \item Um álcool primário poderia ser desidratado pelo mesmo mecanismo?
\end{enumerate}

\medskip
\noindent\textbf{Parte IV --- Massas.}
\begin{enumerate}[resume]
  \item \qty{10.0}{g} de 2-metilpropan-2-ol são convertidos no
        \omterm{def:b1:alcohol-activation:alkoxide}{alcóxido}. Calcule a quantidade de matéria.
  \item Que massa de iodometano é necessária para um excesso de \qty{10}{\percent}?
  \item Calcule a massa teórica de MTBE.
  \item Calcule a massa obtida para um \omterm{def:b1:extent-q-and-k:fractional-extent}{rendimento} de \qty{75}{\percent}.
\end{enumerate}
\end{problem}
